% 原书物理页=印刷页8--11；1.3 全节.
\section{拓扑性质}\label{sec:lie-1-3}
\begin{numpara}[单位元分支]\label{np:lie-1-3-1}
设 $G$ 为一个 Lie 群（或更一般地，一个拓扑群）. $G$ 中包含单位元 $e$ 的连通分支通常记作 $G^0$，并称为 ``\newterm{lie-identity-component}{单位元分支}{identity component}''.
\end{numpara}
\begin{proposition*}
设 $G$ 为一个 Lie 群（或更一般地，一个局部连通的群）. 那么 $G^0$ 是 $G$ 中既开又闭的\newterm{lie-normal-subgroup}{正规子群}{normal subgroup}，并且由 $e$ 的任意一个连通邻域生成.
\end{proposition*}
\begin{proof}
$G^0$ 与一切连通分支一样是闭的；又因 $G$ 局部连通，它是开的. 映射 $x\mapsto x^{-1}$ 连续且把 $e$ 送到 $e$，故把 $G^0$ 送到 $G^0$ 内. 同样，映射 $(x,y)\mapsto xy$ 把 $G^0\times G^0$ 送到 $G^0$ 内. 由此可见 $G^0$ 是一个群. 沿用同一论证，对任意 $g\in G$，映射 $x\mapsto gxg^{-1}$ 把 $G^0$ 送到 $G^0$ 内，故 $G^0$ 是正规的.
现在设 $\mathcal U$ 是 $e$ 在 $G$ 中的一个连通邻域，令 $\mathcal V:=\bigcup_{n=1}^{\infty}\mathcal U^n$，其中 $\mathcal U^n=\mathcal U\cdot\mathcal U^{n-1}$. 它是一些交非空的连通集之并，因而是 $e$ 的一个连通邻域，故 $\mathcal V\subset G^0$. 令 $\mathcal W:=\mathcal U\cap\mathcal U^{-1}$，它是 $e$ 的一个对称邻域. 那么 $\mathcal V$ 包含 $\bigcup_{n=1}^{\infty}\mathcal W^n$，后者是 $G$ 的一个开子群，因而包含 $G^0$. 既然 $\mathcal V=G^0$，$\mathcal U$ 的确生成 $G^0$.
\end{proof}
于是商群 $G/G^0$ 是一个离散群，通常记作 $\pi_0(G)$，并称为 ``$G$ 的\newterm{lie-component-group}{连通分支群}{group of connected components}''.
在此顺便指出，做一做或重做下面的习题是有益的：
\begin{exercise*}
回顾：若有满射 $\pi:X\twoheadrightarrow Y$，其中 $X$ 为拓扑空间，我们定义 $Y$ 上的\newterm{lie-quotient-topology}{商拓扑}{quotient topology}为使 $\pi$ 连续的 $Y$ 上最细的拓扑. 这使得我们能够在 $H$ 为拓扑群 $G$ 的子群时，给商空间 $G/H$ 赋予拓扑. 试验证：
\begin{itemize}[label=--]
\item $\pi$ 是开映射.
\item $G/H$ 是分离的当且仅当 $H$ 在 $G$ 中闭.
\item $G/H$ 是离散的当且仅当 $H$ 在 $G$ 中开.
\end{itemize}
\end{exercise*}
正如我们已经指出的，Lie 群这一类包含离散群类，特别地包含有限群类. 微分技巧对这些群的研究显然无所助益，因而通常只能给出所研究 Lie 群的单位元分支的信息. 例如，上面的命题蕴含：
\begin{corollary*}
设 $\varphi:G\longrightarrow H$ 为一个光滑同态，$H$ 连通. 若 $d_e\varphi:T_eG\longrightarrow T_eH$ 是满射，则 $\varphi$ 是满射.
\end{corollary*}
\begin{proof}
由秩定理，$\operatorname{Im}(\varphi)$ 包含 $e$ 的一个邻域. 由于它是一个群且 $H$ 连通，上面的命题给出 $\operatorname{Im}(\varphi)=H$.
\end{proof}
回顾一下，其逆命题也成立：$\varphi$ 满 $\Rightarrow d_e\varphi$ 满，参见注~\ref{rem:lie-1-1-10}.
\begin{numpara}[经典群的连通性]\label{np:lie-1-3-2}
\textit{线性群.}\ $\mathrm{GL}_n(\mathbb C)$、$\mathrm{SL}_n(\mathbb C)$ 与 $\mathrm{SL}_n(\mathbb R)$ 都是连通的. 相反，$\mathrm{GL}_n(\mathbb R)^0=\{g\in\mathrm{GL}_n(\mathbb R),\ \det(g)>0\}$，而 $\pi_0(\mathrm{GL}_n(\mathbb R))=\mathbb Z/2\mathbb Z$.
\textit{正交群.}\ 对 $K=\mathbb R$ 或 $\mathbb C$，有 $\mathrm{O}(n,K)^0=\mathrm{SO}(n,K)$，且 $\pi_0(\mathrm{O}(n,K))=\mathbb Z/2\mathbb Z$. 在 $K=\mathbb R$ 上且当 $pq\neq0$ 时，有 $\pi_0(\mathrm{SO}(p,q))=\mathbb Z/2\mathbb Z$ 与 $\pi_0(\mathrm{O}(p,q))\simeq\mathbb Z/2\mathbb Z\times\mathbb Z/2\mathbb Z$.
\textit{酉群与辛群.}\ 它们全都连通.
下面这个习题用来证明 $\mathrm{SO}(n,\mathbb R)$ 的连通性.
\begin{exercise*}
考虑 $\mathrm{SO}(n,\mathbb R)$ 在球面 $S^{n-1}\subset\mathbb R^n$ 上的自然作用.
\begin{itemize}[label=--]
\item 验证向量 $(0,\cdots,0,1)$ 的稳定子群可等同于 $\mathrm{SO}(n-1,\mathbb R)$.
\item 证明由此得到的映射 $\mathrm{SO}(n,\mathbb R)/\mathrm{SO}(n-1,\mathbb R)\longrightarrow S^{1}$ootnote{译注：原文此处印作 $S^{1}$，按上下文应为 $S^{n-1}$；照录原文.} 是一个同胚（利用 $\mathrm{SO}(n,\mathbb R)$ 与 $S^{n-1}$ 的紧性）.
\item 证明：若 $H$ 是 $G$ 的一个闭的连通子群，则商映射诱导双射 $\pi_0(G)\simeq\pi_0(G/H)$. 由此用归纳法推出 $\mathrm{SO}(n,\mathbb R)$ 的连通性.
\end{itemize}
\end{exercise*}
对 $\mathrm{U}(n)$（作用在 $S^{2n-1}$ 上）或 $\mathrm{GL}_n$（作用在 $K^n$ 上）可以如法炮制. 在后一情形，第三点因缺乏紧性而更为细致，参见 M1 讲义 Thm 1.1.10.
\end{numpara}
\begin{numpara}[单连通性]\label{np:lie-1-3-3}
回顾代数拓扑中的几个定义. 设 $X$ 为一个拓扑空间.
\begin{itemize}[label=--]
\item $X$ 的一个\newterm{lie-trivial-covering}{平凡覆盖}{trivial covering}是形如 $Y=X\times I\xrightarrow{\ pr_1\ }X$ 的映射，其中 $I$ 是一个离散集合.
\item $X$ 的一个\newterm{lie-covering}{覆盖}{covering}（局部平凡的）是映射 $Y\xrightarrow{\ \pi\ }X$，使得每一点 $x\in X$ 都有邻域 $\mathcal U$，使 $\pi|_{\pi^{-1}(\mathcal U)}$ 是 $\mathcal U$ 的一个平凡覆盖.
\item 若 $X$ 连通，则当 $X$ 的每个覆盖都是平凡的时，称 $X$ 为\newterm{lie-simply-connected}{单连通}{simply connected}的.
\item 若 $X$ \newterm{lie-path-connected}{道路连通}{path-connected}且局部单连通（例如一个\newterm{lie-topological-manifold}{拓扑流形}{topological manifold}），则存在一个覆盖 $\widetilde X\xrightarrow{\ \pi\ }X$，其源 $\widetilde X$ 单连通. 人们姑且称之为 $X$ 的\newterm{lie-universal-covering}{万有覆盖}{universal covering}，尽管它仅在不唯一地相差一个同构的意义下才是唯一的（参见下文）. 群 $\operatorname{Aut}_X(\widetilde X)$（因而仅在相差同构的意义下有定义）就是 $X$ 的\newterm{lie-fundamental-group}{基本群}{fundamental group}.
\item 固定一点 $x\in X$，并选取 $x$ 的一个提升 $\widetilde x\in\widetilde X$. 带基点的覆盖 $(\widetilde X,\widetilde x)\longrightarrow(X,x)$ 便具有如下泛性质：对任何带基点的覆盖 $(Y,y)\xrightarrow{\ p\ }(X,x)$，存在唯一的带基点覆盖 $(\widetilde X,\widetilde x)\xrightarrow{\ q\ }(Y,y)$ 使得 $\pi=p\circ q$. 于是偶对 $(\widetilde X,\widetilde x)$ 在相差唯一同构的意义下唯一.
\item 由泛性质，映射 $\operatorname{Aut}_X(\widetilde X)\longrightarrow\pi^{-1}(x)$，$\varphi\mapsto\varphi(\widetilde x)$ 是双射，从而在 $\pi^{-1}(x)$ 上赋予群结构. 如此得到的群记作 $\pi_1(X,x)$，它在相差唯一同构的意义下有良好定义.
\end{itemize}
当 $X$ 是光滑流形时，$X$ 的任何覆盖都典范地是光滑流形之间的光滑映射. 特别地，$X$ 的万有覆盖是一个光滑流形.
\noindent\textbf{例.}
\begin{itemize}[label=--]
\item 对非零整数 $n$，映射 $\mathbb C^{\times}\longrightarrow\mathbb C^{\times}$，$z\mapsto z^n$ 是一个覆盖.
\item 指数映射 $\mathbb C\xrightarrow{\ \exp\ }\mathbb C^{\times}$ 是 $\mathbb C^{\times}$ 的一个万有覆盖.
\item 若有限群 $\Gamma$ 自由地作用在一个拓扑空间上，则商映射 $\pi:X\longrightarrow X/\Gamma$ 是一个覆盖（例如：$X=\mathrm{SL}_n(\mathbb C)$，$\Gamma=Z(\mathrm{SL}_n(\mathbb C))=\mu_n$，$X/\Gamma=\mathrm{PSL}_n(\mathbb C)$）. 若 $X$ 单连通，则 $\pi_1(X/\Gamma)\simeq\Gamma$.
\item 若 $X$ 与 $Y$ 都单连通，则 $X\times Y$ 也单连通.
\end{itemize}
\end{numpara}
\begin{remark*}
基本群 $\pi_1(X,x)$ 也可以定义为以 $x$ 为基点的闭路的同伦类之集，并配以由闭路的拼接所诱导的运算. 读者可将其作为习题加以验证，或查阅一本代数拓扑教材.
\end{remark*}
\begin{proposition}\label{prop:lie-1-3-4}
设 $G$ 为一个连通 Lie 群，$\widetilde G\xrightarrow{\ \pi\ }G$ 为 $G$ 的一个万有覆盖. 在纤维 $\pi^{-1}(e)$ 中取定一个元素 $\widetilde e$. 那么在 $\widetilde G$ 上存在唯一的 Lie 群结构，它以 $\widetilde e$ 为单位元，并使 $\pi$ 成为 Lie 群同态. 此外，$\operatorname{Ker}(\pi)\simeq\pi_1(G)$ 是 $\widetilde G$ 的一个离散的中心子群.
\end{proposition}
\begin{proof}
在 $e$ 的纤维中取定一个元素 $\widetilde e\in\pi^{-1}(e)$，并注意 $\widetilde G\times\widetilde G$ 是单连通的. 于是 $m$ 有唯一的提升 $\widetilde m$，即唯一一个使图表
\[
\begin{tikzcd}
\widetilde G\times\widetilde G \ar[r,"{\widetilde m}"] & \widetilde G, \\
G\times G \ar[u,"{\pi\times\pi}"'] \ar[r,"m"] & G \ar[u,"\pi"']
\end{tikzcd}
\]
交换的连续映射 $\widetilde m$，它把 $(\widetilde e,\widetilde e)$ 送到 $\widetilde e$. 我们观察到，$\widetilde m\circ(\mathrm{id}_{\widetilde G}\times\widetilde m):\widetilde G\times\widetilde G\times\widetilde G\longrightarrow\widetilde G$ 是 $m\circ(\mathrm{id}_G\times m)$ 的一个提升，它把 $(\widetilde e,\widetilde e,\widetilde e)$ 送到 $\widetilde e$. 同样，$\widetilde m\circ(\widetilde m\times\mathrm{id}_{\widetilde G})$ 是 $m\circ(m\times\mathrm{id}_G)$ 的一个提升，它也把 $(\widetilde e,\widetilde e,\widetilde e)$ 送到 $\widetilde e$. 而由于 $m$ 是结合的，我们有 $m\circ(m\times\mathrm{id}_G)=m\circ(\mathrm{id}_G\times m)$. 由带基点提升的唯一性，由此推得 $\widetilde m\circ(\mathrm{id}_{\widetilde G}\times\widetilde m)=\widetilde m\circ(\widetilde m\times\mathrm{id}_{\widetilde G})$，故 $\widetilde m$ 是 $\widetilde G$ 上一个结合的运算.
同样，逆映射 $i:G\longrightarrow G$ 有唯一的提升 $\widetilde i:\widetilde G\longrightarrow\widetilde G$，它把 $\widetilde e$ 送到 $\widetilde e$.``$i$ 是与 $m$ 相伴的逆映射''这一事实等价于等式 $m\circ(\mathrm{id}_G\times i)=\mathrm{id}_G$. 映射 $\widetilde m\circ(\mathrm{id}_{\widetilde G}\times\widetilde i)$ 是 $m\circ(\mathrm{id}_G\times i)$ 的提升，且把 $\widetilde e$ 送到 $\widetilde e$，正如映射 $\mathrm{id}_{\widetilde G}$ 那样. 提升的唯一性于是蕴含 $m\circ(\mathrm{id}_G\times i)=\mathrm{id}_{\widetilde G}$\footnote{译注：原文如此；按上下文此处应为 $\widetilde m\circ(\mathrm{id}_{\widetilde G}\times\widetilde i)=\mathrm{id}_{\widetilde G}$，等号左端才是 $\widetilde G$ 上的映射.}，从而 $\widetilde i$ 是 $\widetilde m$ 的``逆''映射，故 $\widetilde m$ 是一个群运算.
由构造，$\pi$ 是满射，其核在 $\widetilde G$ 中离散且正规. 由于 $\widetilde G$ 连通，它通过共轭作用在 $\operatorname{Ker}\pi$ 上的作用必是平凡的，即 $\operatorname{Ker}\pi$ 是中心的. 如上面的第三个例子那样，且由于 $G=\widetilde G/\operatorname{Ker}\pi$，我们有 $\operatorname{Ker}\pi\simeq\pi_1(G)$.
\end{proof}
\begin{numpara}[经典群的基本群]\label{np:lie-1-3-5}
若 $H\subset G$ 是连通 Lie 群且 $H$ 闭，我们稍后将证明商映射 $G\twoheadrightarrow G/H$ 是以 $H$ 为纤维的拓扑纤维化（甚至是光滑纤维化）. 在这些条件下，与纤维化相伴的长正合列（参见课程或代数拓扑教材）给出正合列
\[
\pi_1(H)\longrightarrow\pi_1(G)\longrightarrow\pi_1(G/H)\longrightarrow\{1\}
\]
（此处诸基本群都取在单位元处）. 这使得我们能用与计算 $\pi_0$ 相同的策略（在球面上的作用等）来计算群 $\mathrm{SU}(n)$ 与 $\mathrm{SO}(n)$ 的 $\pi_1$. 利用极分解（参见 M1 讲义 1.2.4），可以过渡到线性群. 以下就是所得结果.
\textit{酉群.}\ 群 $\mathrm{SU}(n)$ 单连通，而行列式表明 $\pi_1(\mathrm{U}(n))\simeq\pi_1(\mathrm{U}(1))\simeq\mathbb Z$.
\textit{正交群.}\ $\pi_1(\mathrm{SO}(2))\simeq\mathbb Z$（容易）. 一个经典习题（见下文）证明 $\mathrm{SU}(2)$ 是 $\mathrm{SO}(3)$ 的两叶覆盖，这蕴含 $\pi_1(\mathrm{SO}(3))\simeq\mathbb Z/2\mathbb Z$. 由此归纳地推出，当 $n\geqslant3$ 时 $\pi_1(\mathrm{SO}(n))=\mathbb Z/2\mathbb Z$，但需要在上述正合列的左端补上一项 $\pi_2(G/H)$，并且要知道当 $n\geqslant3$ 时 $\pi_2(S^n)=0$.$\mathrm{SO}(n)$ 的万有覆盖记作 $\mathrm{Spin}(n)$（\newterm{lie-spin-group}{旋量群}{spin group}）. 例如 $\mathrm{Spin}(3)\simeq\mathrm{SU}(2)$.
\textit{线性群.}\ 群 $\mathrm{SL}_n(\mathbb C)$ 单连通，而行列式给出 $\pi_1(\mathrm{GL}_n(\mathbb C))\simeq\pi_1(\mathbb C^{\times})\simeq\mathbb Z$. 相反，$\pi_1(\mathrm{SL}_n(\mathbb R))\simeq\pi_1(\mathrm{GL}_n(\mathbb R)^0)=\mathbb Z/2\mathbb Z$. 我们稍后将证明 $\mathrm{SL}_n(\mathbb R)$ 的万有覆盖是一个\emph{非线性} Lie 群，即不能嵌入任何 $\mathrm{GL}_m(\mathbb R)$ 的 Lie 群.
\textit{辛群.}\ $\mathrm{Sp}_{2n}(\mathbb C)$ 单连通，但 $\pi_1(\mathrm{Sp}_{2n}(\mathbb R))\simeq\mathbb Z$. 特别地，$\mathrm{Sp}_{2n}(\mathbb R)$ 具有唯一的连通两叶覆盖. 人们称之为 ``\newterm{lie-metaplectic-group}{metaplectic 群}{metaplectic group}''；它也是一个非线性 Lie 群的例子.
\end{numpara}
